\notepage{66}{A vector identity for Pauli matrices}
The identity
\[
 \boxed{(\mathbf A\cdot\boldsymbol\sigma)(\mathbf B\cdot\boldsymbol\sigma)
 = (\mathbf A\cdot\mathbf B)I
   +\ii(\mathbf A\times\mathbf B)\cdot\boldsymbol\sigma}
\]
is useful, for example, in the relativistic quantum mechanics of electrons
and the Dirac equation. Here the components of $\mathbf A$ and $\mathbf B$
are numbers, or commute with each other and with the Pauli matrices.
\smallheading{Derivation}
Expanding the product gives
\begin{align*}
 (\mathbf A\cdot\boldsymbol\sigma)(\mathbf B\cdot\boldsymbol\sigma)
 &=\left(\sum_i A_i\sigma_i\right)\left(\sum_j B_j\sigma_j\right)\\
 &=\sum_{i,j}A_iB_j\sigma_i\sigma_j\\
 &=\sum_i A_iB_i\sigma_i^2
   +\sum_{(i,j)\in\mathcal C}A_iB_j\sigma_i\sigma_j
   +\sum_{(i,j)\in\mathcal C}A_jB_i\sigma_j\sigma_i,
\end{align*}
where $\mathcal C=\{(1,2),(2,3),(3,1)\}$. The diagonal part is
$(\mathbf A\cdot\mathbf B)I$. Pairing the off-diagonal terms,
\begin{align*}
 \sum_{i\ne j}A_iB_j\sigma_i\sigma_j
 &=\sum_{(i,j)\in\mathcal C}
    \left(A_iB_j\sigma_i\sigma_j+A_jB_i\sigma_j\sigma_i\right)\\
 &=\sum_{(i,j)\in\mathcal C}(A_iB_j-A_jB_i)\sigma_i\sigma_j\\
 &=\ii\bigl[(A_2B_3-A_3B_2)\sigma_1
 +(A_3B_1-A_1B_3)\sigma_2\\
 &\hspace{39mm}+(A_1B_2-A_2B_1)\sigma_3\bigr]\\
 &=\ii(\mathbf A\times\mathbf B)\cdot\boldsymbol\sigma.
\end{align*}
